\documentclass[11pt]{article}
%% Language and font encodings
\usepackage[english]{babel}
\usepackage[utf8x]{inputenc}
\usepackage[T1]{fontenc}

\usepackage{helvet}

%% Sets page size and margins
\usepackage[letterpaper,top=3cm,bottom=2cm,left=3cm,right=3cm,marginparwidth=1.75cm]{geometry}

%% Useful packages
\usepackage{amsmath}
\usepackage{graphicx}
\usepackage{tcolorbox}
\usepackage{amssymb}
\usepackage{amsthm}
\usepackage{lastpage}
\usepackage{accents}
\usepackage{multicol}
\usepackage{booktabs}

% For better list numbering
\usepackage[shortlabels]{enumitem}

% Font
% \usepackage{tgbonum}


% Tikz
\usepackage{tikz}

\usetikzlibrary{calc,fit,shapes.misc,backgrounds}
\usepackage{pgfplots}
\pgfplotsset{compat = newest}
\usetikzlibrary{positioning, arrows.meta}
\usepgfplotslibrary{fillbetween}

% Headers
\usepackage{fancyhdr}
\pagestyle{fancy}

% Store \@title as \thetitle
\makeatletter
\let\thetitle\@title
\makeatother

\fancyhf{}
\lhead{\fontfamily{qbk}\fontsize{10}{11}\selectfont Name:}
\rhead{\fontfamily{qbk}\fontsize{10}{11}\selectfont ECON 3070-020}
\rfoot{\fontfamily{qbk}\fontsize{10}{11}\selectfont \thepage}


% Sections and Subsections

% define colors
\definecolor{buff-gold}{HTML}{CFB87C}
\definecolor{buff-grey}{HTML}{565A5C}
% custom tcolorbox
\tcbset{colframe=buff-gold, colback=white!100!black}

% new page per section
\usepackage{titlesec}
\newcommand{\sectionbreak}{\clearpage}
% change style of section
\usepackage{sectsty}
\sectionfont{\color{buff-gold} \fontfamily{qbk}\selectfont}
\subsectionfont{\color{buff-grey} \fontfamily{qbk}\selectfont}


\begin{document}
  
\section*{Problem Set 3}

\emph{Each student must submit your own \textbf{HANDWRITTEN} work. Turn in a paper copy at the beginning of lecture on 10/5/2026.  Please write your name on each page and make sure your problem set is stapled.  Students are allowed to work in groups, but each student must turn in their own set of answers.  No help from outside of the course is allowed.  To receive full credit for a problem, all work must be shown.  Partial credit will be awarded.  All problems are weighted evenly.}


\begin{enumerate}
  \item For each of the production functions below, answer the following questions:
  \begin{enumerate}[(i)]
    \item What is the marginal product of each of the inputs?
    
    \item Does the marginal product of $L$ diminish, remain constant, or increase as the level of $L$ increases?
    
    \item What is the marginal rate of substitution ($MRTS$) of $L$ and $K$?
    
    \item Is the $MRTS_{L,K}$ diminishing, constant, or increasing as the firm substitutes more $L$ for $K$, holding the level of output constant?
    
    \item Does the production function exhibit increasing, constant, or decreasing returns to scale? Show your work.
  \end{enumerate}

  \begin{enumerate}[(a)]
    \item $Q(K, L) = 6 L^{\frac{1}{2}} K^{\frac{1}{2}}$

\newpage     
    \item $Q(K, L) = \sqrt{L} + K$
    
    \vspace*{100mm}
    \item $Q(K, L) = K + 2L$

\newpage
	\item $Q(K, L) = L^{3}K^{3}$
  \end{enumerate}

  \newpage
  \item Suppose that a firm originally has the production function $Q(K, L) = 10L + 10K$. Over time as the company learns, the production function changes to $Q(K, L) = 40L + 20K$.
  
  \begin{enumerate}[(a)]
    \item Show that the innovation has resulted in technological progress in the sense defined in the notes.
    
    \vspace*{40mm}
    \item Is the technological progress neutral, labor-saving, or capital-saving? How can you tell?
  \end{enumerate}
  
  \vspace*{40mm}
  \item Suppose that a firm's production function is given by $Q = KL +K$.  At point A, the firm uses $K= 3$ units of capital and $L = 5$ units of labor.  At point B, along the same isoquant, the firm uses on $K = 1$ units of capital.  First calculate how much labor is required at point B.  Then calculate the elastitcity of substitution between points A and B.

\newpage
  \item For each of the following parts draw a series of isoquants such that $Q_{1} < Q_{2} < Q_{3}$ that are appropriate for the given elasticity of substitution.  Give an example of a production function that has an elasticity of substitution of the given value.
  
  \begin{enumerate}[(a)]
    \item $\sigma = \infty$
    
    
    \vspace*{60mm}
    \item $\sigma = 0$
    
    
    \vspace*{60mm}
    \item $\sigma = 1$
  \end{enumerate}


  \newpage
  \item Let $T$ represent car tires and $F$ represent car frames. The production of a car requires 4 tires and 1 frame. 
  \begin{enumerate}[(a)]
    \item Draw at least two isoquants for car production, with tires on the horizontal axis and frames on the vertical axis. Be sure to label the axes and the isoquants.
    
    \vspace*{100mm}
    \item Write the production function in mathematical notation, $Q(T, F)$, for the car producer's production function.
  \end{enumerate}
\end{enumerate}


\end{document}